% =========================================================
% Computational Linear Algebra --- Lengths and Dot Products
% (Strang \S{}1.2 / Johnston \S{}1.2)
% =========================================================

\subsection{Lengths and Dot Products}

% ---------------------------------------------------------
\subsection*{The Dot Product}

\begin{definitionbox}{Definition (Dot Product)}
\begin{definition}[Dot Product]
\label{def:dot-product}
The \emph{dot product} (or \emph{inner product}) of $\mathbf{v}, \mathbf{w} \in \mathbb{R}^n$ is
\[
  \mathbf{v} \cdot \mathbf{w} = v_1 w_1 + v_2 w_2 + \cdots + v_n w_n \;\in\; \mathbb{R}.
\]
\end{definition}
\end{definitionbox}
\begin{remark*}[Standard quantified statement]
This definition fixes the named object, operation, or relation by the formula
or conditions stated in the definition.
\end{remark*}

\begin{remark*}[Interpretation]
The statement records the local meaning of the named concept used by the
surrounding computational or geometric construction.
\end{remark*}

\DefinitionalRoot

\begin{remark*}[Exposition]
\textbf{Properties of the Dot Product.}
For all $\mathbf{v}, \mathbf{w}, \mathbf{x} \in \mathbb{R}^n$ and $c \in \mathbb{R}$:
\begin{enumerate}
  \item $\mathbf{v} \cdot \mathbf{w} = \mathbf{w} \cdot \mathbf{v}$ \hfill (commutativity)
  \item $\mathbf{v} \cdot (\mathbf{w} + \mathbf{x}) = \mathbf{v} \cdot \mathbf{w} + \mathbf{v} \cdot \mathbf{x}$ \hfill (distributivity)
  \item $\mathbf{v} \cdot (c\mathbf{w}) = c(\mathbf{v} \cdot \mathbf{w})$ \hfill (scalar compatibility)
  \item $\mathbf{v} \cdot \mathbf{v} \geq 0$, with equality iff $\mathbf{v} = \mathbf{0}$ \hfill (positive definiteness)
\end{enumerate}
\end{remark*}

% ---------------------------------------------------------
\subsection*{Length and Angle}

\begin{definitionbox}{Definition (Length and Unit Vector)}
\begin{definition}[Length and Unit Vector]
\label{def:length-unit-vector}
The \emph{length} (or \emph{norm}) of $\mathbf{v} \in \mathbb{R}^n$ is
\[
  \|\mathbf{v}\| = \sqrt{\mathbf{v} \cdot \mathbf{v}} = \sqrt{v_1^2 + v_2^2 + \cdots + v_n^2}.
\]
A \emph{unit vector} has length $1$.
The \emph{normalisation} of a nonzero vector $\mathbf{v}$ is
$\hat{\mathbf{v}} = \mathbf{v}/\|\mathbf{v}\|$, which is the unit vector in the direction of $\mathbf{v}$.
\end{definition}
\end{definitionbox}
\begin{remark*}[Standard quantified statement]
This definition fixes the named object, operation, or relation by the formula
or conditions stated in the definition.
\end{remark*}

\begin{remark*}[Interpretation]
The statement records the local meaning of the named concept used by the
surrounding computational or geometric construction.
\end{remark*}

\DefinitionalRoot

\begin{remark*}[Exposition]
\textbf{Cauchy--Schwarz Inequality.}
For all $\mathbf{v}, \mathbf{w} \in \mathbb{R}^n$:
\[
  |\mathbf{v} \cdot \mathbf{w}| \leq \|\mathbf{v}\|\,\|\mathbf{w}\|,
\]
with equality if and only if $\mathbf{v}$ and $\mathbf{w}$ are parallel
(one is a scalar multiple of the other).
\end{remark*}

\begin{definitionbox}{Definition (Angle Between Vectors)}
\begin{definition}[Angle Between Vectors]
\label{def:angle-between-vectors}
The \emph{angle} $\theta \in [0, \pi]$ between nonzero vectors $\mathbf{v}, \mathbf{w} \in \mathbb{R}^n$
is defined by
\[
  \cos\theta = \frac{\mathbf{v} \cdot \mathbf{w}}{\|\mathbf{v}\|\,\|\mathbf{w}\|}.
\]
The vectors are \emph{orthogonal} (written $\mathbf{v} \perp \mathbf{w}$) if $\theta = \pi/2$,
equivalently if $\mathbf{v} \cdot \mathbf{w} = 0$.
\end{definition}
\end{definitionbox}
\begin{remark*}[Standard quantified statement]
This definition fixes the named object, operation, or relation by the formula
or conditions stated in the definition.
\end{remark*}

\begin{remark*}[Interpretation]
The statement records the local meaning of the named concept used by the
surrounding computational or geometric construction.
\end{remark*}

\DefinitionalRoot

\begin{remark*}[Exposition]
\textbf{Dot product and normals in NURBS.}
The unit normal to a parametric surface $S(u,v)$ is computed as
$N = S_u \times S_v / \|S_u \times S_v\|$, which requires both the
cross product and normalisation via the norm.
The dot product also enters the convex hull argument: showing that a
B-spline curve lies on a specific side of a hyperplane reduces to
checking the sign of dot products of control points with the hyperplane's normal.
\end{remark*}
